\chapter{Evaluation Metrics}
\label{app:evaluation_metrics}
\hspace{\parindent}
This appendix defines the main evaluation metrics used throughout this dissertation to analyze
planning performance, regression quality, and the distribution of experimental results. The
definitions are included to support the interpretation of the quantitative results presented in the
evaluation chapters.

\section{Mean Absolute Error}
\label{app:mae}
\hspace{\parindent}
The MAE measures the average absolute difference between an observed value
and the corresponding estimated or predicted value. Given a set of \(n\) observations, where
\(y_i\) denotes the observed value and \(\hat{y}_i\) denotes the estimated value, MAE is defined as

\begin{equation}
    \mathrm{MAE} = \frac{1}{n} \sum_{i=1}^{n} \left| y_i - \hat{y}_i \right|.
\end{equation}

This metric is expressed in the same unit as the evaluated variable and provides a direct measure
of the average magnitude of the error. In this dissertation, MAE is used when the objective is to
quantify prediction or estimation error in an interpretable way, without distinguishing between
positive and negative deviations.

\section{Coefficient of Determination}
\label{app:r_squared}
\hspace{\parindent}
The coefficient of determination, commonly denoted as \(R^2\), measures how much of the variance in
the observed data is explained by a regression model. It is defined as

\begin{equation}
    R^2 = 1 - \frac{\sum_{i=1}^{n} \left(y_i - \hat{y}_i\right)^2}
    {\sum_{i=1}^{n} \left(y_i - \bar{y}\right)^2},
\end{equation}
where \(y_i\) is the observed value, \(\hat{y}_i\) is the predicted value, and \(\bar{y}\) is the
mean of the observed values.

An \(R^2\) value close to 1 indicates that the model explains a large proportion of the observed
variance, while values close to 0 indicate limited explanatory power. Negative values may occur
when the model performs worse than simply predicting the mean of the observed data. In this
dissertation, \(R^2\) is used to assess the explanatory quality of regression models used for
runtime analysis.

\section{Empirical Cumulative Distribution Function}
\label{app:ecdf}
\hspace{\parindent}
The Empirical Cumulative Distribution Function (ECDF) represents the proportion of observations
that are less than or equal to a given value. For a set of \(n\) observations, it can be written as

\begin{equation}
    \hat{F}_n(x) = \frac{k}{n},
\end{equation}
where \(k\) is the number of observations with value less than or equal to \(x\).

In this dissertation, ECDF plots are used to compare experimental runs by showing the percentage of
runs that achieve a value below a given threshold. This is useful for metrics such as planning time
and solution cost, since it shows the distribution of results instead of only average values.